\newpage

\begin{flushright}
  \fechaclase{18 de noviembre de 2026}
\end{flushright}

cuando $e_1=0$

\begin{equation*}
    |e_1|^{\alpha-1}=\frac{1}{|e_1|^{1-\alpha}},\quad 0<\alpha<1
\end{equation*}

lo que ocasiona que

\begin{equation*}
    ||_1|^{\alpha-1}\to \infty
\end{equation*}

por lo cual, la superficie anterior genera una singularidad en la $e_1=0$ proponemos entonces la superficie

\begin{equation*}
    s_T=e_1+\frac{1}{m_T}<e_2>^\gamma, \quad |<\gamma<2
\end{equation*}

derivando la superficie

\begin{equation*}
    \begin{aligned}
        S_T=&\dot{e}_1+\frac{\gamma}{m_T}|e_2|^{\gamma-1}\dot{e}_2\\
            =& e_2+\frac{\gamma}{m_T}|e_2|^{\gamma-1}[\ddot{r}-f(x)-u]
    \end{aligned}
\end{equation*}

proponemos la funcion de Lyapunov

\begin{equation*}
    V=\frac{1}{2}s_T^2
\end{equation*}

derivamos

\begin{equation*}
    \begin{aligned}
        \dot{V}=&s_T\dot{s}_T\\
            =&s_T(e_2+\frac{\gamma}{m_T}|e_2|^{\gamma-1}( \underbrace{\ddot{r}-f(x,t)}_F-u))
    \end{aligned}
\end{equation*}

deseamos que 

\begin{equation*}
    \begin{aligned}
        \dot{V}\leq & (-\frac{\gamma (k-F)}{m_T}|e_2|^{\gamma-1}sign(s_T))\\
            \leq & - \frac{\gamma (k-F)}{m_T} |e_2|^{\gamma-1}|s_T|
    \end{aligned}
\end{equation*}

despejando

\begin{equation*}
    \frac{m_T}{\gamma}e_2+|e_2|^{\gamma-1}(\underbrace{\ddot{\gamma}-f(x,t)}_F-u) = - (k-F)|e_2|^{\gamma-1}sign(s_T)
\end{equation*}

dividiendo entre $|e_2|^{\gamma-1}$

\begin{equation*}
    \frac{m_T}{\gamma}\frac{e_2}{|e_2|^{\gamma-1}}+(F-u)= -(K-F) sign(s_T)
\end{equation*}

\begin{equation*}
    \frac{m_T}{\gamma}<e_2>^{2-\gamma} + (F-u)=-(K-F)sign(s_T)
\end{equation*}

\begin{equation*}
    u=\frac{m_T}{\gamma}<e_2>^{2-\gamma}+k Sign(s_T)
\end{equation*}

mostrando que 

\begin{equation*}
    \dot{V}\leq - \frac{\gamma k}{ m_T}|e_2|^{\gamma-1}|s_T|
\end{equation*}